\printedpage{79}
% head: A. Grothendieck : Dérivateurs, Ch. VII, Édition des Universités de Montpellier II et Paris VII
% foot: Version 22 décembre 2001

This time, the composition functors
\[ C_n(\underline{t};a,b)\times C_n(\underline{t'};b,c) \longrightarrow C_n(\underline{t''};a,c), \]
are compatible with the transition functors for varying $n$, and one finds
\msp{page 126}
\[ C_\infty(\underline{t};a,b)\times C_\infty(\underline{t'};b,c) \longrightarrow C_\infty(\underline{t''};a,c), \]
hence an arrow in Hot
\[ \comp_{\tau_0} : C_\infty\times C_\infty \longrightarrow C_\infty. \]
We must see that it does not depend on the choice of $\tau_0$, but if we have $\tau_0$ and $\tau'_0$, we compare them by $\tau_0\times\tau'_0 \begin{array}{l}\longrightarrow \tau_0\\ \longrightarrow \tau'_0\end{array}$, so it suffices to see that if there exists an arrow $\tau_0\longrightarrow\tau'_0$, then $\comp_{\tau_0} = \comp_{\tau'_0}$. But such an arrow induces arrows $\underline{t}_{\tau_0}\longrightarrow\underline{t}_{\tau'_0}$, $\underline{t'}_{\tau_0}\longrightarrow\underline{t'}_{\tau'_0}$, $\underline{t''}_{\tau_0}\longrightarrow\underline{t''}_{\tau'_0}$, whence a diagram in Cat
\[
\begin{tikzcd}
C_\infty(\underline{t}_{\tau'_0};a,b)\times C_\infty(\underline{t'}_{\tau'_0};b,c) \arrow[r, "{\comp_{\tau'_0}}"] \arrow[d] & C_\infty(\underline{t''}_{\tau'_0};a,c) \arrow[d] \\
C_\infty(\underline{t}_{\tau_0};a,b)\times C_\infty(\underline{t'}_{\tau_0};b,c) \arrow[r, "{\comp_{\tau_0}}"] & C_\infty(\underline{t''}_{\tau_0};a,c),
\end{tikzcd}
\]
which one checks at once is commutative. Now this diagram can be rewritten in Hot as
\[
\begin{tikzcd}
C_\infty(a,b)\times C_\infty(b,c) \arrow[r, "{\comp_{\tau'_0}}"] \arrow[d, "{\mathrm{id}}"] & C_\infty(a,c) \arrow[d, "{\mathrm{id}}"] \\
C_\infty(a,b)\times C_\infty(b,c) \arrow[r, "{\comp_{\tau_0}}"] & C_\infty(a,c),
\end{tikzcd}
\]
\msp{page 127}
whence $\comp_{\tau_0} = \comp_{\tau'_0}$.

I shrink from the real difficulty - that of proving the associativity of this composition law. It becomes associative in $\Hot^{\wedge}_{0}$ - but for two arrows in Hot to be equal (here then
\[ C_\infty(a,b)\times C_\infty(b,c)\times C_\infty(c,d) \rightrightarrows C_\infty(a,d)\quad ), \]
is it enough that this be so for their images in $\Hot^{\wedge}_{0}$ - in other words, is the functor
\[ \Hot\longrightarrow\Hot^{\wedge}_{0}, \]
(besides being conservative) at least faithful (if not fully faithful)?

\clearpage
\printedpage{80}
% head: A. Grothendieck : Dérivateurs, Ch. VII, Édition des Universités de Montpellier II et Paris VII
% foot: Version 22 décembre 2001

There is also the fact that this composition law is compatible with that of the $H(a,b)$ in Ind\,Hot, for the canonical arrows
\[ H(a,b)\longrightarrow C_\infty(a,b) \qquad \text{in Ind\,Hot}, \]
if we would know that these are epimorphisms, and remain so under cartesian products, we would have won.

Of course, we will know that it is associative, once we know that the composititon defined here is compatible with the composition of paths in the semi-simplicial framework. But how does one proceed in this framework, to compose paths of unbounded length? Is there not exactly the same difficulty here? Unless one gets % [sic] composititon -> composition, misspelling % Reading resolved against the supplied PDF; see corrections.csv.
\msp{page 128}
out of it by reducing at once, for the theory of $\Omega$ operations, to Kan complexes, and to paths of type $\Delta_+$. But to compose them \textit{functorially,} there still seems to me to be the same problem - one will find paths of type $\Delta_+\Delta_+ = \boldsymbol{\Delta}_2^+$, and not $\Delta_+$, but one can then compose the direct homotopies (amounts to restricting to $\Delta_+\simeq\{0,2\}\subseteq\boldsymbol{\Delta}_2^+$). Now I worked out the formalism of categories of paths without worrying about conditions of Kan type. I presume that in the semi-simplicial framework, the $\uHom(\boldsymbol{\Delta}_1,X)$ do \textit{not} give the right path object, for $X$ not Kan. (For otherwise, it would seem that for $X=X_0$ whose components are \textit{finite,} the calculation of the $\pi_i$ ought to be finite, which is nevertheless not the case, I have always been assured \dots)

I presume nevertheless, when one takes $\Sigma=\mathrm{Fl}(X)$, that the objects $\Ch_\infty(X;a,b)$ and $\Ch_\infty(\bar{X};a,b)$ are isomorphic in Hot, when one embeds $X$ in a Thomason-Kan envelope $\bar{X}$.

\clearpage
\printedpage{81}
% head: A. Grothendieck : Dérivateurs, Ch. VII, Édition des Universités de Montpellier II et Paris VII
% foot: Version 22 décembre 2001

\msp{page 129}

\section*{11\quad Quick return to the $(\mathcal{M},\Sigma)$, from the perspective of derivators}

When one is interested in the formalism of the $\Ch_\Sigma(X;a,b)$ with a view to a \textit{derivator} $\mathbf{D}_{X,\Sigma}$ that one would like to define in terms of $X$, $\Sigma$, it would seem that the prederivator one does indeed obtain only deserves the name of derivator when one assumes that for every arrow
\[ Y\longrightarrow Y' \]
in Cat which is a Hot-equivalence, the corresponding arrows
\[ \mathrm{Hom}_{\mathrm{Ind\,Cat\,hot}}(Y',H(a,b)) \longrightarrow \mathrm{Hom}_{\mathrm{Ind\,Cat\,hot}}(Y,H(a,b)), \]
or again
\[ \varinjlim_n \mathrm{Hom}_{\mathrm{Cat\,hot}}(Y',\Ch_{n,\Sigma}(a,b)) \longrightarrow \varinjlim_n \mathrm{Hom}_{\mathrm{Cat\,hot}}(Y,\Ch_{n,\Sigma}(a,b)) \]
are bijections. I looked in V at what this means for the object $H(a,b) = (\Ch_{n,\Sigma}(a,b))_n$ of Ind\,Cat\,hot. It is this condition alone that ensures, from the point of view of derivators, a convincing interpretation of $\Ch_{\infty,\Sigma}(a,b)$,
\[ \Ch_{\infty,\Sigma}(a,b) \;=\; \varinjlim_n \Ch_{n,\Sigma}(a,b) \]
as an object of Hot, or more exactly of $\Hot^{\wedge}_{0}$, for then one has for $Y$ a \textit{finite} category
\[ \underbrace{\mathrm{Hom}_{\mathrm{Cat\,hot}}(Y,H(a,b))}_{\simeq\,\mathrm{Hom}_{X(Y)\Sigma_Y^{-1}}(a_X,b_X)} \;\simeq\; \mathrm{Hom}_{\Hot}(Y,\Ch_{\infty,\Sigma}(a,b)), \]
\msp{page 130}
for I do want the type in Hot to be the right one, that it does correspond to the homotopy of path spaces. So one would like a morphism
\[ X\xrightarrow{\;f\;}X' \]
in (Cat), which is a Hot-equivalence, to define
\[ \Ch_\infty(X;a,b) \longrightarrow \Ch_\infty(X';f(a),f(b)), \]
which are themselves also Hot-equivalences. On the other hand, this is obviously compatible with the brilliant composition of paths, which I managed to define between these objects of Hot. So (subject to having checked what I have just said) the verification of associativity would reduce to the case where $X$ is Thomason-Kan. But in this case, there is a chance that, just as in the semi-simplicial case, the $\Ch_\infty(X;a,b)$ can be defined (as objects of (Hot)) as being simply the paths of type $\Delta_+$ from $a$ to $b$, which one then composes in the trivial way. But no, these categories are discrete! % [sic] ramenerait -> ramènerait, missing accent

I give up for the moment pursuing the question of associativity any further. Ask \textsc{Thomason} about the properties of the functor $\Hot\longrightarrow\Hot^{\wedge}_{0}$.

\clearpage
\printedpage{82}
% head: A. Grothendieck : Dérivateurs, Ch. VII, Édition des Universités de Montpellier II et Paris VII
% foot: Version 22 décembre 2001

\msp{page 131}
This means, very nearly, that one has an interpretation of $\Ch_{\infty,\Sigma}(a,b)$, as an object of Hot, as an ``external Hom'' $\mathrm{Hom\,ext}(a,b)$ with values in Hot, defined on the category $\mathcal{M}=X$ using just the admissible set of arrows $\Sigma\subseteq\mathrm{Fl}(\mathcal{M})$, an operation satisfying the condition
\[ \boxed{\mathrm{Hom}_{\Hot}(Y,\mathrm{Hom\,ext}(a,b)) \;\simeq\; \mathrm{Hom}_{\mathcal{M}(Y)\Sigma_Y^{-1}}(a_X,b_X)} \tag*{$(*)$} \]
(which can also be interpreted, either as
\[ \mathrm{Hom}_{\mathcal{M}}(p_!(a_Y),b) \]
if $p_! : \underbrace{\mathcal{M}(Y)\Sigma_Y^{-1}}_{\mathbf{D}(Y)} \longrightarrow \mathcal{M}$ exists, or as
\[ \mathrm{Hom}_{\mathcal{M}}(a,p_*(b_Y)) \]
if $p_* : \underbrace{\mathcal{M}(Y)\Sigma_Y^{-1}}_{\mathbf{D}(Y)} \longrightarrow \underbrace{\mathcal{M}}_{=\,\mathbf{D}(*)}$ exists); with the one small nuisance that this formula is valid only for $Y$ \textit{finite.} If one wants to express the right-hand side for $Y$ varying arbitrarily in Cat, an object of Hot ($\mathrm{Hom\,ext}(a,b)$) on the left-hand side does not suffice, but one needs the ind-object $H(a,b) = (\Ch_n(a,b))_{n\in\mathbf{N}}$ in Hot. So if one insists on working in a ``domain'' Diag of the derivator $\mathbf{D}=\mathbf{D}_{\mathcal{M},\Sigma}$ as large as possible, while restricting oneself to finite categories,
\msp{page 132}
and one wants a formula $(*)$ to remain valid for every $Y$ in the domain, one must ``in all rigor'' define the $\mathrm{Hom\,ext}(a,b)$ not as an object of Hot, but of Ind\,Hot. It would therefore then be described as
\[ \mathrm{Hom\,ext}(a,b) = (\Ch_{n,\Sigma}(a,b))_n \qquad \text{(ind-object in Hot)}\ , \]
and the object $\Ch_{\infty,\Sigma}$ in Hot is a sort of direct limit of it.

I wondered whether the preliminary condition, namely that the functors $H(a,b)$ on Cat turn Hot-equivalences $Y'\longrightarrow Y$ into bijections, was really a (strong) restriction. I convinced myself that this is indeed the case, already in the special case where $\Sigma=\mathrm{Fl}(X)$. It seems to me that when the connected components of $\Omega(X,a)$ are not contractible, the said condition is practically never satisfied. I thought e.g. of the case where $X$ has the homotopy type (in Hot) of a 2-sphere, so that its $\Omega(X,a)=\Ch_\infty(X,a,a)$ must have a $\pi_1\simeq\mathbf{Z}$. Let us take $Y$ having the type in Hot of a circle, say
\[
\begin{tikzcd}
 & & a & & \\
Y\colon & 0 \arrow[ur] \arrow[dr] & & 0' \arrow[ul] \arrow[dl] & , \\
 & & b & &
\end{tikzcd}
\]

\clearpage
\printedpage{83}
% head: A. Grothendieck : Dérivateurs, Ch. VII, Édition des Universités de Montpellier II et Paris VII
% foot: Version 22 décembre 2001

\msp{page 133}
so the functors
\[ f : Y\longrightarrow \Ch_\infty(X;\bullet,\bullet) \overset{\text{def}}{=} \Omega(X,\bullet) \]
correspond to \textit{certain} ``loops'' in $\Omega(X,\bullet)$ (but loops with unspecified origin), they have a homotopy type in the topological sense (``number of turns'' around the sphere $X$) given by an integer $n\in\mathbf{Z}$ (the ``degree'' of the loop $\lambda$). The obvious intuition is that the set of degrees $\delta$ which one can obtain in this way, with a given \textit{finite} category $Y$ which represents the circle $S^1\in(\Hot)$, is finite - $|\delta|$ surely bounded above by the number of these arrows (hence at least by 4 in the case under consideration, each arrow of $Y$ making one go \textit{at most} once around the pseudo-circle $\Omega(S^2,\bullet)$) (\footnote[28]{probably one finds at most only the degrees $\delta$ between $-2$ and $+2$.}). On the other hand, when one represents $S^1$ by arbitrarily large categories $Y_n$ (say obtained from the intervals $\Phi_n=(\Delta_+\Delta_-)^n$ by identifying source and target), one must find arbitrarily high degrees. Now every $Y_n$ maps to $Y(=Y_2)$ by an arrow
\[ Y_n = Y' \longrightarrow Y = Y_2 \]
\msp{page 134}
which is of degree 1, hence which is a Hot-equivalence. But the corresponding map
\[ \pi_0\uHom(Y,\Omega(X,\bullet)) \longrightarrow \pi_0\uHom(Y',\Omega(X,\bullet)) \]
is not bijective. This kind of argument seems to me bound to work whenever the set of homotopy groups $\pi_j(X,x)$, for $j\geq 2$ (hence the set of homotopy groups $\pi_i$ of $\Omega(X,x)$, for $i\geq 1$) is \textit{infinite,} taking for $Y$, $Y'$ different representatives $Y_n$ of the sphere $S^i$, say the $Y_n$ obtained from $\Phi_n$ $(n\geq 2)$ by taking $\Phi_n^i = \Phi_n\times\Phi_n\times\cdots\times\Phi_n$ (which is a sort of combinatorial $i$-cube), and identifying its ``boundary'' with a single point. I presume that the number of elements in $\pi_i(\Omega(X,x))$ that one can obtain from arrows $Y_n\longrightarrow\Omega(X,x)$ is limited by the size of $Y_n$, but one can surely obtain each of them, by taking $n$ large enough. So if $\xi\in\pi_i(\Omega)$ escapes $Y_{n_0}$ ($n_0=0$, say), one takes an $n'$ large enough that $\xi$ does not escape $Y_{n'}$, and one chooses % [sic] d'élements -> d'éléments, missing accent
\[ Y_{n'} \longrightarrow Y_n\ , \]
\texttt{[}$n=n_0$\texttt{?]}

\msp{page 135}
which is a Hot-equivalence (e.g. with the aid of a retraction of $\Phi_{n'}$ onto $\Phi_n$, compatible with sources and targets).

This draws strong attention to the fact that if $Y$, $C$ are two categories, the category $\uHom(Y,C)$ \textit{is not} a representative of $\uHom_{\Hot}(\mathrm{hot}(Y),\mathrm{hot}(C))$, one has only a canonical arrow
\[ \mathrm{hot}(\uHom(Y,C)) \longrightarrow \uHom_{\Hot}(\mathrm{hot}(Y),\mathrm{hot}(C)), \tag*{$(*)$} \]